\chapter{Replication as a Fixed Point}
\label{sec:replication}
% ===================================================================

\section{The Concurrent Y Combinator}

The Rho calculus admits a \emph{concurrent Y combinator}: a
process $\mathbf{Y}$ satisfying the fixed-point equation
\[
  \mathbf{Y}(F) \;\bisim\; F\bigl(\mathbf{Y}(F)\bigr) \;\mid\; \mathbf{Y}(F).
\]
This is the abstract core of the $\pi$-calculus replication
operator $!P \equiv P \mid !P$, expressed entirely within the
Rho calculus using only the comm rule and quoting, without
adding replication as a new primitive.

\begin{remark}[Replication is Picked Out as a Fixed Point]
The concurrent Y combinator is not just \emph{a} fixed point of
the interaction dynamics; it is the canonical fixed point of the
\emph{replication map} $F \mapsto F \mid F$.
In this sense replication is not an accidental feature of the
Rho calculus but a structural inevitability: any sufficiently
expressive interactive GSLT must contain a fixed point of its own
duplication map, and that fixed point \emph{is} replication.
\end{remark}

\section{The Comm Rule as a Crossover Operator}

The single base rewrite rule of the Rho calculus is:
\[
  \mathtt{for}(y \leftarrow x)\,P \;\mid\; x!(Q)
  \;\rewrite\;
  P\{@Q / y\}
\]
This rule substitutes the \emph{quoted code} $@Q$ --- a name
encoding the entire program $Q$ --- into the body $P$ of the
receiver.
The result is a new process that has incorporated material from
both $P$ (the receiver's structure) and $Q$ (the sender's content).

\begin{remark}[One Step from Genetic Recombination]
A \emph{crossover operator} in the sense of genetic algorithms
recombines fragments of two parent programs to produce an offspring.
The comm rule is one definitional step from this:
instead of substituting $Q$ wholesale into $P$, a crossover
variant would substitute a \emph{fragment} of $Q$ --- a sub-process
selected by some recombination policy.

This places replication, recombination, and variation --- the three
pillars of Darwinian evolution --- all within reach of a single
rule.
Replication is the Y combinator fixed point.
Recombination is the comm rule with fragment substitution.
Variation is the accumulation of errors in the substitution process,
weighted by the amplitude structure of the path integral.
Darwinian evolution is not added to the framework; it is derivable
from it.
\end{remark}

% ===================================================================
